\section{Why econirl} \label{sec:why}

We design \pkg{econirl} for three target audiences.

The first audience is structural econometricians who fit dynamic discrete choice models with parametric rewards and want a path to neural reward representations without leaving their inference toolkit. They have always relied on standard errors, identification diagnostics, and likelihood ratio tests to defend their estimates. They want the same machinery available when the reward is replaced by a neural network.

The second audience is inverse reinforcement learning researchers who train adversarial or distribution-matching reward learners and want standard errors, identification diagnostics, and counterfactual welfare elasticity for the rewards they recover. They have always reported policy quality and reward correlations. They want the inferential vocabulary of structural estimation available without rewriting their training loops.

The third audience is applied empiricists who fit both families on the same panel without reimplementing either. They want to compare a nested fixed point estimator and an adversarial inverse reinforcement learning estimator on a single dataset with the same data loader, the same feature mapping, the same standard error methodology, and the same counterfactual simulator.

\pkg{econirl} differentiates itself along three axes. The first is the unified result object. Every estimator returns an \code{EstimationSummary} that exposes point estimates, standard errors, identification diagnostics, the value function, the policy, and a \code{summary()} method. The shared object means the post-estimation pipeline does not branch on the estimator. The second is the inference layer. Standard errors are available for any estimator that exposes a Hessian, by asymptotic, robust sandwich, nonparametric bootstrap, and clustered methods. The same Hessian feeds the identification diagnostics, the profile likelihood, and the Wald tests. The third is the JAX backbone. Automatic differentiation eliminates hand-coded analytic derivatives, which is the most common source of bugs in structural estimation code. Just-in-time compilation accelerates the inner loops by an order of magnitude relative to interpreted \proglang{NumPy}. Vectorized mapping enables bootstrap resampling and Monte Carlo replications to run in parallel on a single machine without manual thread management.

The contribution of this paper is the design and implementation of these three differentiators. Section~\ref{sec:methods} lays out the notation and the twelve estimators. Section~\ref{sec:design} describes the \code{Estimator} protocol and the inference layer. Section~\ref{sec:illustrations} demonstrates that the unified workflow holds on three real-data panels and across all twelve estimators in a single benchmark.
